\chapter{Guía de lectura}

\SeccionLibro{Objetivos del libro}

\subsection{Objetivo general}

Desarrollar competencias de pensamiento computacional, razonamiento lógico y programación que permitan analizar, modelar, construir, evaluar y mejorar soluciones a problemas mediante procedimientos computacionales.

\subsection{Objetivos específicos}

\begin{itemize}
\item Comprender los fundamentos del pensamiento computacional y su importancia en la resolución de problemas.
\item Desarrollar habilidades de descomposición, reconocimiento de patrones, abstracción y generalización.
\item Identificar los requerimientos, entradas, procesos y salidas de un problema computacional.
\item Diseñar algoritmos mediante pseudocódigo y diagramas de flujo.
\item Implementar soluciones computacionales utilizando los fundamentos del lenguaje Python.
\item Aplicar estructuras secuenciales, selectivas y repetitivas.
\item Organizar soluciones mediante funciones y colecciones de datos.
\item Aplicar técnicas de verificación, validación, depuración y optimización de programas.
\item Integrar los conocimientos adquiridos mediante proyectos prácticos.
\end{itemize}

\SeccionLibro{Público objetivo y alcance}

El presente libro está dirigido principalmente a estudiantes universitarios de los primeros ciclos de Ingeniería Informática, Ingeniería de Sistemas, Ingeniería de Software, Ciencias de la Computación y carreras afines.

También puede ser utilizado por docentes, estudiantes de institutos tecnológicos y personas interesadas en iniciarse de manera autodidacta en el pensamiento computacional y la programación.

No se requieren conocimientos previos avanzados de informática. El contenido parte de conceptos elementales y progresa hacia la construcción de soluciones computacionales de mayor complejidad.

El alcance comprende desde la identificación y el análisis de problemas hasta la implementación, evaluación y mejora de programas. La inteligencia artificial y la visión artificial se presentan como aplicaciones complementarias de estos principios.

\SeccionLibro{Organización y metodología de aprendizaje}

El libro se organiza en dos partes y cuatro capítulos. La \RefTabla{tab:estructura-del-libro} presenta sus contenidos.

\begin{tablalibro}{2}{L{0.29}L{0.71}}
\caption{Estructura del libro}\label{tab:estructura-del-libro}\\
\toprule
\rowcolor{Papel}
\CabeceraTabla{Parte y capítulo} & \CabeceraTabla{Contenido} \\
\midrule
\endfirsthead
\SigueTabla{2}
\toprule
\rowcolor{Papel}
\CabeceraTabla{Parte y capítulo} & \CabeceraTabla{Contenido} \\
\midrule
\endhead
\midrule
\ContinuaTabla{2}
\endfoot
\bottomrule
\endlastfoot

Parte I · \hyperref[cap:01-pensamiento-computacional]{Capítulo~\ref*{cap:01-pensamiento-computacional}} & Fundamentos del pensamiento computacional y análisis de problemas. \\
Parte I · \hyperref[cap:02-diseno-algoritmico]{Capítulo~\ref*{cap:02-diseno-algoritmico}} & Modelado y diseño de soluciones algorítmicas. \\
Parte II · \hyperref[cap:03-fundamentos-de-programacion]{Capítulo~\ref*{cap:03-fundamentos-de-programacion}} & Construcción de soluciones mediante fundamentos de programación. \\
Parte II · \hyperref[cap:04-pruebas-y-mejora]{Capítulo~\ref*{cap:04-pruebas-y-mejora}} & Pruebas, depuración y mejora de soluciones computacionales. \\
\end{tablalibro}


La metodología de aprendizaje se fundamenta en la integración de conocimientos conceptuales y experiencias prácticas. Cada capítulo incorpora explicaciones teóricas, ejemplos desarrollados, representaciones gráficas, casos prácticos, preguntas de comprensión, ejercicios y desafíos de resolución de problemas.

\SeccionLibro{Requisitos previos y herramientas necesarias}

Para desarrollar las actividades propuestas se recomienda disponer de conocimientos básicos de matemática, comprensión lectora y manejo elemental de computadoras.

\begin{tablalibro}{2}{L{0.31}L{0.69}}
\caption{Herramientas para las actividades}\label{tab:herramientas-para-las-actividades}\\
\toprule
\rowcolor{Papel}
\CabeceraTabla{Herramienta} & \CabeceraTabla{Finalidad} \\
\midrule
\endfirsthead
\SigueTabla{2}
\toprule
\rowcolor{Papel}
\CabeceraTabla{Herramienta} & \CabeceraTabla{Finalidad} \\
\midrule
\endhead
\midrule
\ContinuaTabla{2}
\endfoot
\bottomrule
\endlastfoot

Python 3 & Implementación de soluciones computacionales. \\
Visual Studio Code & Escritura y edición de programas. \\
Google Colab & Ejecución de código desde el navegador. \\
PSeInt & Representación y ejecución de pseudocódigo. \\
Draw.io & Elaboración de diagramas de flujo. \\
Jupyter Notebook & Experimentación y documentación de soluciones. \\
GitHub & Organización y almacenamiento de recursos y códigos. \\
\end{tablalibro}


Puede elegir las herramientas según la actividad y los recursos disponibles.

\SeccionLibro{Cómo utilizar el libro y acceder a los recursos}

Recorra los capítulos en orden. Antes de consultar un ejemplo resuelto, formule su propia estrategia y luego compare ambos procedimientos.

Ejecute los archivos de la carpeta \Codigo{recursos/} con Python 3. Varíe las entradas y registre los resultados en las plantillas del \RefAnexo{anexo:plantillas}. Use las preguntas para revisar conceptos y los desafíos para integrar lo aprendido. Puede desarrollar el proyecto transversal de forma individual o en equipo.

\textbf{Convención numérica.} La prosa y las tablas usan coma decimal (14,2). El código Python y sus salidas usan punto decimal (14.40).
